% --- Archon physics formalization source begin ---
% archon:physics
% archon:covers PhyXMiniProblems/problem_phyx_mini_0658.lean
% archon:source-report reports/phyx_mini/problem_phyx_mini_0658.source.json

\chapter{Physics problem phyx\_mini\_0658}
\label{ch:PhyXMiniProblems_problem_phyx_mini_0658}

\paragraph{Problem source.}
\#\# Physical scenario

In most metals, the atomic ions form a regular arrangement called a crystal lattice. The conduction electrons in the sea of electrons move through this lattice. Figure is a one-dimensional model of a crystal lattice. The ions have mass \$m\$, charge \$e\$, and an equilibrium separation \$b\$. Suppose this crystal consists of aluminum ions with an equilibrium spacing of \$0.30\$ \$nm\$.

\#\# Question

What are the energies of the four lowest vibrational states of these ions?

\#\# Answer choices

- A: A: 0.0428 eV

- B: B: 0.0531 eV

- C: C: 0.0634 eV

- D: D: 0.0832 eV

\#\# Figure caption (auxiliary; use the image as primary evidence)

The image shows three circles, each with a red plus sign inside, indicating positive charges. These circles are aligned horizontally and evenly spaced.

- Two double-headed arrows are present:
  - The top arrow is positioned above the middle circle, oriented horizontally, suggesting directionality or distance between the outer charges.
  - The lower arrows are placed between each pair of circles, also horizontal, pointing toward or away from the central circle.

- The symbol "b" is located below each of the lower arrows, indicating the distance between each adjacent pair of charges.

The arrangement suggests a symmetric distribution of three positive charges, each separated by a distance "b".

\#\# Dataset metadata

Quantum Phenomena; Numerical Reasoning, Physical Model Grounding Reasoning

\paragraph{Recorded answer/context.}
C: C: 0.0634 eV

\paragraph{Figure/image path.}
/root/proposal\_for\_physic/science-mango/phyx\_mini\_run/phyx\_data/test\_image/658.png

\paragraph{Formalization target.}
create a compiling Lean file with sorry bodies at `PhyXMiniProblems/problem\_phyx\_mini\_0658.lean`. The Lean declarations must preserve the physical quantities, dimensions or dimensional roles, figure labels, governing-law hypotheses, and final relation expressed by this problem.
Use Mathlib/Physlib names found through LeanExplore where available. If a physics API is missing, introduce faithful local abstractions rather than scalar placeholder aliases.

\begin{theorem}[Physics formalization target]
\label{thm:physics:phyx_mini_0658:target}
\lean{PhyXMiniProblems.ProblemPhyXMini0658.aluminum_ion_four_lowest_vibrational_energies}
\uses{def:physics:phyx-mini-0658:phyxminiproblems-problemphyxmini0658-energyinjoules, def:physics:phyx-mini-0658:phyxminiproblems-problemphyxmini0658-energyinelectronvolts, def:physics:phyx-mini-0658:phyxminiproblems-problemphyxmini0658-aluminumionlatticesetup, def:physics:phyx-mini-0658:phyxminiproblems-problemphyxmini0658-matchesaluminumcrystalscenario, def:physics:phyx-mini-0658:phyxminiproblems-problemphyxmini0658-matchesprimarylatticefigure, def:physics:phyx-mini-0658:phyxminiproblems-problemphyxmini0658-usesroundedaluminumreferencedata, def:physics:phyx-mini-0658:phyxminiproblems-problemphyxmini0658-satisfiestwoneighborcoulomblaw, def:physics:phyx-mini-0658:phyxminiproblems-problemphyxmini0658-realizessmalloscillationmodel, def:physics:phyx-mini-0658:phyxminiproblems-problemphyxmini0658-satisfiesvibrationalschrodingerequation, def:physics:phyx-mini-0658:phyxminiproblems-problemphyxmini0658-energyquantuminelectronvolts, def:physics:phyx-mini-0658:phyxminiproblems-problemphyxmini0658-answerchoice}
The assigned autoformalize agent should translate this physics problem into Lean declarations in the covered file, with theorem and lemma proof bodies written as `by sorry`.
\end{theorem}
\begin{proof}
This is an autoformalization task, not a proof task. Produce statements that can later be proved by the physics prover without weakening the physical model.
\end{proof}

% --- Archon declaration topology begin ---
\section{Declaration topology}
\begin{definition}[spring Constant Dimension]
  \label{def:physics:phyx-mini-0658:phyxminiproblems-problemphyxmini0658-springconstantdimension}
  \lean{PhyXMiniProblems.ProblemPhyXMini0658.springConstantDimension}
  The dimension of a spring constant, N/m = kg/s²
\end{definition}
\begin{proof}
  This is the declared model definition of spring Constant Dimension.
\end{proof}
\begin{definition}[Length Quantity]
  \label{def:physics:phyx-mini-0658:phyxminiproblems-problemphyxmini0658-lengthquantity}
  \lean{PhyXMiniProblems.ProblemPhyXMini0658.LengthQuantity}
  A nonnegative, unit-independent physical length
\end{definition}
\begin{proof}
  This is the declared model definition of Length Quantity.
\end{proof}
\begin{definition}[Mass Quantity]
  \label{def:physics:phyx-mini-0658:phyxminiproblems-problemphyxmini0658-massquantity}
  \lean{PhyXMiniProblems.ProblemPhyXMini0658.MassQuantity}
  A nonnegative, unit-independent physical mass
\end{definition}
\begin{proof}
  This is the declared model definition of Mass Quantity.
\end{proof}
\begin{definition}[Charge Magnitude Quantity]
  \label{def:physics:phyx-mini-0658:phyxminiproblems-problemphyxmini0658-chargemagnitudequantity}
  \lean{PhyXMiniProblems.ProblemPhyXMini0658.ChargeMagnitudeQuantity}
  A nonnegative, unit-independent magnitude of electric charge
\end{definition}
\begin{proof}
  This is the declared model definition of Charge Magnitude Quantity.
\end{proof}
\begin{definition}[Spring Constant Quantity]
  \label{def:physics:phyx-mini-0658:phyxminiproblems-problemphyxmini0658-springconstantquantity}
  \lean{PhyXMiniProblems.ProblemPhyXMini0658.SpringConstantQuantity}
  \uses{def:physics:phyx-mini-0658:phyxminiproblems-problemphyxmini0658-springconstantdimension}
  A nonnegative, unit-independent effective spring constant
\end{definition}
\begin{proof}
  This is the declared model definition of Spring Constant Quantity.
\end{proof}
\begin{definition}[length Readout]
  \label{def:physics:phyx-mini-0658:phyxminiproblems-problemphyxmini0658-lengthreadout}
  \lean{PhyXMiniProblems.ProblemPhyXMini0658.lengthReadout}
  \uses{def:physics:phyx-mini-0658:phyxminiproblems-problemphyxmini0658-lengthquantity}
  Read a physical length in a selected Physlib length unit
\end{definition}
\begin{proof}
  This is the declared model definition of length Readout.
\end{proof}
\begin{definition}[length In Meters]
  \label{def:physics:phyx-mini-0658:phyxminiproblems-problemphyxmini0658-lengthinmeters}
  \lean{PhyXMiniProblems.ProblemPhyXMini0658.lengthInMeters}
  \uses{def:physics:phyx-mini-0658:phyxminiproblems-problemphyxmini0658-lengthquantity, def:physics:phyx-mini-0658:phyxminiproblems-problemphyxmini0658-lengthreadout}
  Coherent-SI metre readout of a physical length
\end{definition}
\begin{proof}
  This is the declared model definition of length In Meters.
\end{proof}
\begin{definition}[length In Nanometers]
  \label{def:physics:phyx-mini-0658:phyxminiproblems-problemphyxmini0658-lengthinnanometers}
  \lean{PhyXMiniProblems.ProblemPhyXMini0658.lengthInNanometers}
  \uses{def:physics:phyx-mini-0658:phyxminiproblems-problemphyxmini0658-lengthquantity, def:physics:phyx-mini-0658:phyxminiproblems-problemphyxmini0658-lengthreadout}
  Nanometre readout of a physical length
\end{definition}
\begin{proof}
  This is the declared model definition of length In Nanometers.
\end{proof}
\begin{definition}[mass In Kilograms]
  \label{def:physics:phyx-mini-0658:phyxminiproblems-problemphyxmini0658-massinkilograms}
  \lean{PhyXMiniProblems.ProblemPhyXMini0658.massInKilograms}
  \uses{def:physics:phyx-mini-0658:phyxminiproblems-problemphyxmini0658-massquantity}
  Coherent-SI kilogram readout of a physical mass
\end{definition}
\begin{proof}
  This is the declared model definition of mass In Kilograms.
\end{proof}
\begin{definition}[charge Readout]
  \label{def:physics:phyx-mini-0658:phyxminiproblems-problemphyxmini0658-chargereadout}
  \lean{PhyXMiniProblems.ProblemPhyXMini0658.chargeReadout}
  \uses{def:physics:phyx-mini-0658:phyxminiproblems-problemphyxmini0658-chargemagnitudequantity}
  Read a charge magnitude in a selected Physlib charge unit
\end{definition}
\begin{proof}
  This is the declared model definition of charge Readout.
\end{proof}
\begin{definition}[charge In Coulombs]
  \label{def:physics:phyx-mini-0658:phyxminiproblems-problemphyxmini0658-chargeincoulombs}
  \lean{PhyXMiniProblems.ProblemPhyXMini0658.chargeInCoulombs}
  \uses{def:physics:phyx-mini-0658:phyxminiproblems-problemphyxmini0658-chargemagnitudequantity, def:physics:phyx-mini-0658:phyxminiproblems-problemphyxmini0658-chargereadout}
  Coherent-SI coulomb readout of a physical charge magnitude
\end{definition}
\begin{proof}
  This is the declared model definition of charge In Coulombs.
\end{proof}
\begin{definition}[charge In Elementary Charges]
  \label{def:physics:phyx-mini-0658:phyxminiproblems-problemphyxmini0658-chargeinelementarycharges}
  \lean{PhyXMiniProblems.ProblemPhyXMini0658.chargeInElementaryCharges}
  \uses{def:physics:phyx-mini-0658:phyxminiproblems-problemphyxmini0658-chargemagnitudequantity, def:physics:phyx-mini-0658:phyxminiproblems-problemphyxmini0658-chargereadout}
  Read a charge magnitude as a multiple of the elementary charge e
\end{definition}
\begin{proof}
  This is the declared model definition of charge In Elementary Charges.
\end{proof}
\begin{definition}[spring Constant In Newtons Per Meter]
  \label{def:physics:phyx-mini-0658:phyxminiproblems-problemphyxmini0658-springconstantinnewtonspermeter}
  \lean{PhyXMiniProblems.ProblemPhyXMini0658.springConstantInNewtonsPerMeter}
  \uses{def:physics:phyx-mini-0658:phyxminiproblems-problemphyxmini0658-springconstantquantity}
  Coherent-SI N/m readout of an effective spring constant
\end{definition}
\begin{proof}
  This is the declared model definition of spring Constant In Newtons Per Meter.
\end{proof}
\begin{definition}[energy In Joules]
  \label{def:physics:phyx-mini-0658:phyxminiproblems-problemphyxmini0658-energyinjoules}
  \lean{PhyXMiniProblems.ProblemPhyXMini0658.energyInJoules}
  Coherent-SI joule readout of a physical energy
\end{definition}
\begin{proof}
  This is the declared model definition of energy In Joules.
\end{proof}
\begin{definition}[electron Volt In Joules]
  \label{def:physics:phyx-mini-0658:phyxminiproblems-problemphyxmini0658-electronvoltinjoules}
  \lean{PhyXMiniProblems.ProblemPhyXMini0658.electronVoltInJoules}
  \uses{def:physics:phyx-mini-0658:phyxminiproblems-problemphyxmini0658-energyinjoules}
  The joule value of Physlib's dimensionful electron volt
\end{definition}
\begin{proof}
  This is the declared model definition of electron Volt In Joules.
\end{proof}
\begin{definition}[energy In Electron Volts]
  \label{def:physics:phyx-mini-0658:phyxminiproblems-problemphyxmini0658-energyinelectronvolts}
  \lean{PhyXMiniProblems.ProblemPhyXMini0658.energyInElectronVolts}
  \uses{def:physics:phyx-mini-0658:phyxminiproblems-problemphyxmini0658-energyinjoules, def:physics:phyx-mini-0658:phyxminiproblems-problemphyxmini0658-electronvoltinjoules}
  Electron-volt readout of a physical energy
\end{definition}
\begin{proof}
  This is the declared model definition of energy In Electron Volts.
\end{proof}
\begin{definition}[Ion Site]
  \label{def:physics:phyx-mini-0658:phyxminiproblems-problemphyxmini0658-ionsite}
  \lean{PhyXMiniProblems.ProblemPhyXMini0658.IonSite}
  The three positive ions displayed in the primary figure
\end{definition}
\begin{proof}
  This is the declared model definition of Ion Site.
\end{proof}
\begin{definition}[Adjacent Gap]
  \label{def:physics:phyx-mini-0658:phyxminiproblems-problemphyxmini0658-adjacentgap}
  \lean{PhyXMiniProblems.ProblemPhyXMini0658.AdjacentGap}
  The two adjacent gaps whose distances are both labelled b
\end{definition}
\begin{proof}
  This is the declared model definition of Adjacent Gap.
\end{proof}
\begin{definition}[Figure Distance Label]
  \label{def:physics:phyx-mini-0658:phyxminiproblems-problemphyxmini0658-figuredistancelabel}
  \lean{PhyXMiniProblems.ProblemPhyXMini0658.FigureDistanceLabel}
  The literal symbolic distance label visible under each gap arrow
\end{definition}
\begin{proof}
  This is the declared model definition of Figure Distance Label.
\end{proof}
\begin{definition}[Ion Species]
  \label{def:physics:phyx-mini-0658:phyxminiproblems-problemphyxmini0658-ionspecies}
  \lean{PhyXMiniProblems.ProblemPhyXMini0658.IonSpecies}
  Atomic species needed by the physical scenario
\end{definition}
\begin{proof}
  This is the declared model definition of Ion Species.
\end{proof}
\begin{definition}[Crystal Lattice Figure]
  \label{def:physics:phyx-mini-0658:phyxminiproblems-problemphyxmini0658-crystallatticefigure}
  \lean{PhyXMiniProblems.ProblemPhyXMini0658.CrystalLatticeFigure}
  \uses{def:physics:phyx-mini-0658:phyxminiproblems-problemphyxmini0658-ionsite, def:physics:phyx-mini-0658:phyxminiproblems-problemphyxmini0658-adjacentgap, def:physics:phyx-mini-0658:phyxminiproblems-problemphyxmini0658-figuredistancelabel}
  Literal presentation data from image 658
\end{definition}
\begin{proof}
  This is the declared model definition of Crystal Lattice Figure.
\end{proof}
\begin{definition}[Aluminum Ion Lattice Setup]
  \label{def:physics:phyx-mini-0658:phyxminiproblems-problemphyxmini0658-aluminumionlatticesetup}
  \lean{PhyXMiniProblems.ProblemPhyXMini0658.AluminumIonLatticeSetup}
  \uses{def:physics:phyx-mini-0658:phyxminiproblems-problemphyxmini0658-lengthquantity, def:physics:phyx-mini-0658:phyxminiproblems-problemphyxmini0658-massquantity, def:physics:phyx-mini-0658:phyxminiproblems-problemphyxmini0658-chargemagnitudequantity, def:physics:phyx-mini-0658:phyxminiproblems-problemphyxmini0658-springconstantquantity, def:physics:phyx-mini-0658:phyxminiproblems-problemphyxmini0658-adjacentgap, def:physics:phyx-mini-0658:phyxminiproblems-problemphyxmini0658-ionspecies, def:physics:phyx-mini-0658:phyxminiproblems-problemphyxmini0658-crystallatticefigure}
  Physical data for the displayed lattice cell and the vibrating middle ion. The map vibrationalEnergy stores the energy observable at each quantum level; it is not defined from any answer choice
\end{definition}
\begin{proof}
  This is the declared model definition of Aluminum Ion Lattice Setup.
\end{proof}
\begin{definition}[Matches Aluminum Crystal Scenario]
  \label{def:physics:phyx-mini-0658:phyxminiproblems-problemphyxmini0658-matchesaluminumcrystalscenario}
  \lean{PhyXMiniProblems.ProblemPhyXMini0658.MatchesAluminumCrystalScenario}
  \uses{def:physics:phyx-mini-0658:phyxminiproblems-problemphyxmini0658-lengthinnanometers, def:physics:phyx-mini-0658:phyxminiproblems-problemphyxmini0658-chargeinelementarycharges, def:physics:phyx-mini-0658:phyxminiproblems-problemphyxmini0658-aluminumionlatticesetup}
  The material and spacing stated in the problem source
\end{definition}
\begin{proof}
  This is the declared model definition of Matches Aluminum Crystal Scenario.
\end{proof}
\begin{definition}[Matches Primary Lattice Figure]
  \label{def:physics:phyx-mini-0658:phyxminiproblems-problemphyxmini0658-matchesprimarylatticefigure}
  \lean{PhyXMiniProblems.ProblemPhyXMini0658.MatchesPrimaryLatticeFigure}
  \uses{def:physics:phyx-mini-0658:phyxminiproblems-problemphyxmini0658-aluminumionlatticesetup}
  All unambiguous geometric and charge evidence read from the primary image. The image supplies no energy value
\end{definition}
\begin{proof}
  This is the declared model definition of Matches Primary Lattice Figure.
\end{proof}
\begin{definition}[Uses Rounded Aluminum Reference Data]
  \label{def:physics:phyx-mini-0658:phyxminiproblems-problemphyxmini0658-usesroundedaluminumreferencedata}
  \lean{PhyXMiniProblems.ProblemPhyXMini0658.UsesRoundedAluminumReferenceData}
  \uses{def:physics:phyx-mini-0658:phyxminiproblems-problemphyxmini0658-massinkilograms, def:physics:phyx-mini-0658:phyxminiproblems-problemphyxmini0658-aluminumionlatticesetup}
  Rounded independent constants used by the textbook numerical estimate. The mass 4.5 × 10⁻²⁶ kg is the usual two-significant-figure aluminum-ion mass, and 9 × 10⁹ N m²/C² is the corresponding rounded Coulomb constant
\end{definition}
\begin{proof}
  This is the declared model definition of Uses Rounded Aluminum Reference Data.
\end{proof}
\begin{definition}[Satisfies Two Neighbor Coulomb Law]
  \label{def:physics:phyx-mini-0658:phyxminiproblems-problemphyxmini0658-satisfiestwoneighborcoulomblaw}
  \lean{PhyXMiniProblems.ProblemPhyXMini0658.SatisfiesTwoNeighborCoulombLaw}
  \uses{def:physics:phyx-mini-0658:phyxminiproblems-problemphyxmini0658-lengthinmeters, def:physics:phyx-mini-0658:phyxminiproblems-problemphyxmini0658-chargeincoulombs, def:physics:phyx-mini-0658:phyxminiproblems-problemphyxmini0658-springconstantinnewtonspermeter, def:physics:phyx-mini-0658:phyxminiproblems-problemphyxmini0658-aluminumionlatticesetup}
  For a signed displacement x of the middle ion, the left neighbor contributes a force toward the right and the right neighbor contributes a force toward the left. Their exact one-dimensional Coulomb-force difference is valid in the open interval |x| < b. The derivative at equilibrium defines the negative effective spring constant. Neither clause contains a vibrational energy or an answer-choice value
\end{definition}
\begin{proof}
  This is the declared model definition of Satisfies Two Neighbor Coulomb Law.
\end{proof}
\begin{definition}[Realizes Small Oscillation Model]
  \label{def:physics:phyx-mini-0658:phyxminiproblems-problemphyxmini0658-realizessmalloscillationmodel}
  \lean{PhyXMiniProblems.ProblemPhyXMini0658.RealizesSmallOscillationModel}
  \uses{def:physics:phyx-mini-0658:phyxminiproblems-problemphyxmini0658-massinkilograms, def:physics:phyx-mini-0658:phyxminiproblems-problemphyxmini0658-springconstantinnewtonspermeter, def:physics:phyx-mini-0658:phyxminiproblems-problemphyxmini0658-aluminumionlatticesetup}
  The dimensionful mass and spring constant are supplied to Physlib's classical oscillator through coherent-SI readouts. Its frequency is then the frequency of Physlib's quantum oscillator for the middle ion
\end{definition}
\begin{proof}
  This is the declared model definition of Realizes Small Oscillation Model.
\end{proof}
\begin{definition}[Satisfies Vibrational Schrodinger Equation]
  \label{def:physics:phyx-mini-0658:phyxminiproblems-problemphyxmini0658-satisfiesvibrationalschrodingerequation}
  \lean{PhyXMiniProblems.ProblemPhyXMini0658.SatisfiesVibrationalSchrodingerEquation}
  \uses{def:physics:phyx-mini-0658:phyxminiproblems-problemphyxmini0658-energyinjoules, def:physics:phyx-mini-0658:phyxminiproblems-problemphyxmini0658-aluminumionlatticesetup}
  The time-independent Schrödinger eigenstate law for the dimensionful energy observable. This governing equation does not assume the closed-form eigenvalue, any of the requested energy readouts, or a multiple-choice option. Physlib independently proves the same equation with quantumOscillator.eigenValue n as its coefficient; comparing the two equations is the later proof route to the energy formula
\end{definition}
\begin{proof}
  This is the declared model definition of Satisfies Vibrational Schrodinger Equation.
\end{proof}
\begin{theorem}[effective Spring Constant eq coulomb Linearization]
  \label{thm:physics:phyx-mini-0658:phyxminiproblems-problemphyxmini0658-effectivespringconstant-eq-coulomblinearization}
  \lean{PhyXMiniProblems.ProblemPhyXMini0658.effectiveSpringConstant_eq_coulombLinearization}
  \uses{def:physics:phyx-mini-0658:phyxminiproblems-problemphyxmini0658-lengthinmeters, def:physics:phyx-mini-0658:phyxminiproblems-problemphyxmini0658-chargeincoulombs, def:physics:phyx-mini-0658:phyxminiproblems-problemphyxmini0658-springconstantinnewtonspermeter, def:physics:phyx-mini-0658:phyxminiproblems-problemphyxmini0658-aluminumionlatticesetup, def:physics:phyx-mini-0658:phyxminiproblems-problemphyxmini0658-matchesaluminumcrystalscenario, def:physics:phyx-mini-0658:phyxminiproblems-problemphyxmini0658-satisfiestwoneighborcoulomblaw}
  The exact effective spring constant obtained by differentiating the two-neighbor Coulomb force at the equilibrium position
\end{theorem}
\begin{proof}
  The conclusion follows from the governing relations and figure readouts named in the statement of effective Spring Constant eq coulomb Linearization.
\end{proof}
\begin{definition}[energy Quantum In Electron Volts]
  \label{def:physics:phyx-mini-0658:phyxminiproblems-problemphyxmini0658-energyquantuminelectronvolts}
  \lean{PhyXMiniProblems.ProblemPhyXMini0658.energyQuantumInElectronVolts}
  \uses{def:physics:phyx-mini-0658:phyxminiproblems-problemphyxmini0658-electronvoltinjoules, def:physics:phyx-mini-0658:phyxminiproblems-problemphyxmini0658-aluminumionlatticesetup}
  The oscillator quantum ℏω, expressed in electron volts
\end{definition}
\begin{proof}
  This is the declared model definition of energy Quantum In Electron Volts.
\end{proof}
\begin{definition}[Answer Choice]
  \label{def:physics:phyx-mini-0658:phyxminiproblems-problemphyxmini0658-answerchoice}
  \lean{PhyXMiniProblems.ProblemPhyXMini0658.AnswerChoice}
  The four answer choices printed in the problem source
\end{definition}
\begin{proof}
  This is the declared model definition of Answer Choice.
\end{proof}
\begin{definition}[answer Choice Energy Electron Volts]
  \label{def:physics:phyx-mini-0658:phyxminiproblems-problemphyxmini0658-answerchoiceenergyelectronvolts}
  \lean{PhyXMiniProblems.ProblemPhyXMini0658.answerChoiceEnergyElectronVolts}
  \uses{def:physics:phyx-mini-0658:phyxminiproblems-problemphyxmini0658-answerchoice}
  The electron-volt number printed beside an answer choice
\end{definition}
\begin{proof}
  This is the declared model definition of answer Choice Energy Electron Volts.
\end{proof}
\begin{definition}[answer Agreement Tolerance Electron Volts]
  \label{def:physics:phyx-mini-0658:phyxminiproblems-problemphyxmini0658-answeragreementtoleranceelectronvolts}
  \lean{PhyXMiniProblems.ProblemPhyXMini0658.answerAgreementToleranceElectronVolts}
  Numerical tolerance appropriate to the rounded textbook input data
\end{definition}
\begin{proof}
  This is the declared model definition of answer Agreement Tolerance Electron Volts.
\end{proof}
\begin{definition}[agrees With Fourth Level]
  \label{def:physics:phyx-mini-0658:phyxminiproblems-problemphyxmini0658-answerchoice-agreeswithfourthlevel}
  \lean{PhyXMiniProblems.ProblemPhyXMini0658.AnswerChoice.agreesWithFourthLevel}
  \uses{def:physics:phyx-mini-0658:phyxminiproblems-problemphyxmini0658-energyinelectronvolts, def:physics:phyx-mini-0658:phyxminiproblems-problemphyxmini0658-aluminumionlatticesetup, def:physics:phyx-mini-0658:phyxminiproblems-problemphyxmini0658-answerchoice, def:physics:phyx-mini-0658:phyxminiproblems-problemphyxmini0658-answerchoiceenergyelectronvolts, def:physics:phyx-mini-0658:phyxminiproblems-problemphyxmini0658-answeragreementtoleranceelectronvolts}
  A generic criterion saying that a printed value agrees, to the stated tolerance, with the fourth-lowest level (n = 3). It does not privilege C
\end{definition}
\begin{proof}
  This is the declared model definition of agrees With Fourth Level.
\end{proof}
% --- Archon declaration topology end ---
% --- Archon physics formalization source end ---
